\begin{tikzpicture}[font=\small]
% (a) point-segment distance and segment-circle intersection
\begin{scope}[scale=0.9]
  \coordinate (A) at (0,0);
  \coordinate (B) at (5,0);
  \coordinate (C) at (2,0.5);
  \coordinate (Q) at (2,0);
  \coordinate (D) at (4,1.6);
  \coordinate (QD) at (4,0);
  % the circle that meets the segment
  \fill[sbRed!12] (C) circle (1);
  \draw[sbRed,thick] (C) circle (1);
  \fill[sbRed] (C) circle (1.3pt);
  \node[sbannot,text=sbRed] at (2,1.1) {$\vect{c}$};
  % the circle that misses it
  \draw[sbGreen,thick] (D) circle (1);
  \fill[sbGreen] (D) circle (1.3pt);
  \node[sbannot,text=sbGreen] at (4.55,1.95) {$\vect{c}'$};
  \draw[sbGreen,dashed] (D) -- (QD);
  \node[sbannot,text=sbGreen] at (4.35,0.8) {$d>r$};
  % the segment
  \draw[sbedge,line width=1.4pt] (A) -- (B);
  \fill[black] (A) circle (1.4pt) node[below=2pt] {$\vect{a}$};
  \fill[black] (B) circle (1.4pt) node[below=2pt] {$\vect{b}$};
  % closest point and distance
  \draw[sbRed,dashed] (C) -- (Q);
  \fill[sbRed] (Q) circle (1.3pt);
  \node[sbannot,text=sbRed] at (2.05,-0.3) {$\vect{q}$};
  \node[sbannot,text=sbRed] at (2.3,0.28) {$d$};
  % entry and exit points of the ray a + t (b - a)
  \fill[sbBlue] (1.134,0) circle (1.4pt);
  \fill[sbBlue] (2.866,0) circle (1.4pt);
  \node[sbannot,text=sbBlue] at (1.1,0.32) {$t_1$};
  \node[sbannot,text=sbBlue] at (2.95,-0.3) {$t_2$};
  \node[sbannot,text=black] at (2.5,-0.95) {(a) segment against two discs};
\end{scope}
% (b) tangent lines from a point to a circle
\begin{scope}[xshift=7.4cm,scale=0.75]
  \coordinate (P) at (0,0);
  \coordinate (C) at (4,0);
  \coordinate (T1) at (3,1.732);
  \coordinate (T2) at (3,-1.732);
  % the cone between the two tangent lines (preview of the collision cone)
  \fill[sbOrange!15] (P) -- (6,3.464) -- (6,-3.464) -- cycle;
  \draw[sbOrange,thick] (P) -- (6,3.464);
  \draw[sbOrange,thick] (P) -- (6,-3.464);
  \fill[sbRed!12] (C) circle (2);
  \draw[sbRed,thick] (C) circle (2);
  \draw[sbGray,dashed] (P) -- (C);
  \draw[sbRed] (C) -- (T1);
  \draw[sbRed] (C) -- (T2);
  \node[sbannot,text=sbRed] at (3.75,1.05) {$r$};
  % right angle at the tangent point
  \draw pic[draw=sbGray,angle radius=2.2mm] {right angle=C--T1--P};
  \fill[black] (P) circle (1.4pt) node[left=2pt] {$\vect{p}$};
  \fill[sbRed] (C) circle (1.4pt) node[below=3pt] {$\vect{c}$};
  \fill[sbOrange] (T1) circle (1.6pt) node[above left=1pt] {$\vect{t}_{+}$};
  \fill[sbOrange] (T2) circle (1.6pt) node[below left=1pt] {$\vect{t}_{-}$};
  % angles: theta at p (half-angle of the cone), alpha at c
  \draw pic[draw=sbGray,angle radius=11mm,"$\theta$",font=\footnotesize,angle eccentricity=1.25] {angle=C--P--T1};
  \draw pic[draw=sbGray,angle radius=7mm,"$\alpha$",font=\footnotesize,angle eccentricity=1.4] {angle=T1--C--P};
  \node[sbannot] at (1.3,1.35) {$\ell=\sqrt{d^2-r^2}$};
  \node[sbannot] at (2,-0.3) {$d$};
  \node[sbannot,text=black] at (3,-3.3) {(b) tangents from $\vect{p}$};
\end{scope}
\end{tikzpicture}
